\chapter*{Abstract}
\addcontentsline{toc}{chapter}{Abstract}

\noindent
Il presente progetto di semestre, svolto in collaborazione con il
Politecnico di Milano, ha per oggetto la \textbf{predizione di parametri
di celle Litio-Ione} caratterizzate tramite spettroscopia di impedenza
elettrochimica galvanostatica (\acrshort{geis}) su una pouch
\acrshort{licoo2} da \SI{10}{\ampere\hour}. Il dataset, fornito dal
gruppo di ricerca del Politecnico, conta
\num{9805} misurazioni organizzate come \textbf{5 livelli di
invecchiamento $\times$ 8 temperature $\times$ 5 \acrshort{soc}} e
$\sim$49 frequenze logaritmicamente spaziate, ovvero \textbf{200 curve di
Nyquist complete}.

Il \textbf{compito primario} consisteva nello sviluppare tre notebook
Jupyter contenenti modelli di machine learning per altrettanti task
predittivi:
\begin{enumerate}[leftmargin=*]
  \item \textbf{Task~1 — Leave-One-Out (LOO):} ricostruzione di una
        singola coppia (Aging, Temperatura) a partire dalle altre 39
        (regressione multi-output).
  \item \textbf{Task~2 — Young/Old Classification:} dati un livello
        di Aging e un livello di \acrshort{soc}, predire la classe
        \emph{young} (Aging $\leq 2$) vs \emph{old} (Aging $\geq 3$)
        delle \textbf{8 curve di Nyquist} (una per temperatura)
        corrispondenti a quella coppia, tenute fuori dall'addestramento.
  \item \textbf{Task~3 — Aging Interpolation:} predizione di un
        \emph{intero} livello di invecchiamento mai osservato in
        addestramento.
\end{enumerate}

I migliori modelli ottenuti su un benchmark di cinque/sei stimatori
sono rispettivamente \emph{KNN distance-weighted} ($R^{2} = 0.991$,
Task 1), \emph{SVM con kernel RBF} (accuracy $= 1.000$ sulle 8 curve
del hold-out di default $(Aging = 4, SOC = 3)$, Task 2) e di nuovo
\emph{KNN distance-weighted} ($R^{2} = 0.986$, Task 3). I risultati
sono prodotti da una pipeline interamente riproducibile (random seed
fissato, splitter coerenti con il task) e tracciata con MLflow.

\medskip
\noindent
\textbf{Estensione volontaria.} Sebbene la consegna fosse limitata ai
notebook, abbiamo costruito \emph{sopra} di essi una piccola piattaforma
MLOps end-to-end -- non richiesta -- per esporre i risultati a un utente
non-developer e per esplorare nella pratica le prassi di settore:
configurazione centralizzata in \acrshort{yaml}, logging strutturato con
rotazione, suite di 81 test pytest con coverage gate al 90\,\% in CI
(attuale: 95\,\%),
scansioni di sicurezza (\texttt{bandit}, \texttt{pip-audit}), build
Docker multistage, backend \texttt{FastAPI} con frontend React per il
\emph{serving} e un modulo di \textbf{drift detection} basato su test di
\acrshort{ks} e \acrshort{psi}.

Il report descrive in dettaglio il dominio fisico, il dataset, la
metodologia condivisa fra i tre task, i risultati ottenuti e -- in
appendice concettuale -- l'architettura della piattaforma volontaria,
chiudendosi con un'analisi dei limiti residui e delle direzioni di
sviluppo futuro.

\chapter*{Abstract (English)}
\addcontentsline{toc}{chapter}{Abstract (English)}

\noindent
This semester project, carried out in collaboration with Politecnico di
Milano, focuses on the \textbf{prediction of Lithium-Ion battery
parameters} characterized via galvanostatic electrochemical impedance
spectroscopy (\acrshort{geis}) on a \SI{10}{\ampere\hour}
\acrshort{licoo2} pouch cell. The dataset, supplied by the Politecnico
research group, contains \num{9805} measurements organized as
\textbf{5 aging levels $\times$ 8 temperatures $\times$ 5 SOC} and
$\sim$49 logarithmically spaced frequencies, i.e.~\textbf{200 complete
Nyquist curves}.

The \textbf{primary deliverable} consisted of three Jupyter notebooks
implementing machine-learning models for three predictive tasks:
\textbf{(1)} Leave-One-Out reconstruction of a single $(Aging,
Temperature)$ pair from the other 39 (multi-output regression);
\textbf{(2)} given an Aging level and a SOC level, predict the
\emph{young} (Aging $\leq 2$) vs \emph{old} (Aging $\geq 3$) class of
the \textbf{8 Nyquist curves} (one per temperature) corresponding to
that pair, held out from training; \textbf{(3)} interpolation of an
entire unseen aging level (multi-output regression).

The best models on a benchmark of five/six estimators are respectively
\emph{KNN distance-weighted} ($R^{2} = 0.991$, Task 1), \emph{SVM with
RBF kernel} (accuracy $= 1.000$ on the 8 curves of the default
hold-out $(Aging = 4, SOC = 3)$, Task 2) and again \emph{KNN
distance-weighted} ($R^{2} = 0.986$, Task 3). Results are produced by
a fully reproducible pipeline (fixed random seed, task-coherent
splitters) and tracked with MLflow.

\medskip
\noindent
\textbf{Voluntary extension.} On top of the notebooks we built a small
end-to-end MLOps platform -- not part of the original deliverable -- to
expose the models to a non-developer user and to gain hands-on
experience with industry MLOps practices: centralized YAML
configuration, structured rotating logging, an 81-test pytest suite
with a 90\,\% coverage gate in CI (currently 95\,\%), security scanning (\texttt{bandit},
\texttt{pip-audit}), multi-stage Docker build, FastAPI backend with
React frontend for serving, and a \textbf{drift detection} module based
on KS and PSI tests.

The report covers the physical domain, the dataset, the methodology
shared across the three tasks, the obtained results and -- as
conceptual appendix -- the architecture of the voluntary platform,
closing with an analysis of residual limitations and directions for
future work.
